\section{两种 Scaling Paradigm 与两条产品研究路径}

讲者先区分 next-token prediction 与 RL on chain of thought，再区分“熟悉界面承载陌生能力”和“先有产品信念，再训练模型做到”。这两个二维坐标非常重要：前一个坐标回答能力怎样获得，后一个坐标回答能力怎样进入产品。把它们混在一起，容易误以为更强预训练会自动给出正确产品行为。

\subsection{Paradigm 1：Next-token prediction 是 world-building machine}

自回归预训练最常见的目标是
\[
\mathcal{L}_{\text{NLL}}(\theta)
=-\sum_{t=1}^{T}\log p_{\theta}(x_t\mid x_{<t}).
\]
$x_t$ 是当前 token，$x_{<t}$ 是历史上下文，$\theta$ 是模型参数。统一目标能让模型吸收语言、知识、风格和部分推理模式，因此被讲者称为 world-building machine；但它奖励的是局部预测，不直接定义一个长任务何时完成，也不提供外部工具成功、用户满意或长期一致性的信号。

\lecturefigure{slide-11-next-token-scaling.jpg}{Paradigm 1：next-token prediction 在多领域建立世界模型}{00:07:58--00:08:39}

\begin{knowledgebox}{为什么早期 token 错误会放大}
生成分布可写为 $p(x_{1:T})=\prod_t p(x_t\mid x_{<t})$。模型一旦采样到偏离目标的 token，后续条件分布就建立在新的历史上；长文本中的角色、事实和计划因此可能逐步漂移。这个问题不是说自回归一定失败，而是说明“局部似然高”与“长期任务成功”不是同一指标。
\end{knowledgebox}

\subsection{Paradigm 2：RL on CoT 把长任务结果写进目标}

强化学习把一次推理或行动轨迹记为 $\tau=(s_0,a_0,\dots,s_T)$，优化
\[
J(\theta)=\mathbb{E}_{\tau\sim\pi_{\theta}}[R(\tau)].
\]
$\pi_\theta$ 是策略，$R(\tau)$ 是整条轨迹的奖励。若 reward 来自代码测试、数学验证器或工具结果，模型就能获得超出逐 token 模仿的任务级信号。实践中常加入参考策略约束，避免策略为追逐奖励过度漂移：
\[
J_{\text{reg}}(\theta)=\mathbb{E}[R(\tau)]
-\beta\,D_{\mathrm{KL}}(\pi_{\theta}\Vert\pi_{\mathrm{ref}}).
\]
$\beta$ 控制行为改进与分布保持的权衡。

\lecturefigure{slide-12-rl-cot-scaling.jpg}{Paradigm 2：在 Chain-of-Thought 轨迹上用 RL 优化复杂任务}{00:08:39--00:09:03}

\begin{warningbox}{CoT 不是自动真实的“内心过程”}
Chain-of-Thought 可以提供更长的计算轨迹、可插入工具调用和中间检查，但可读文字不保证 faithful。模型可能给出正确答案配错误解释，或把奖励偏好的格式学成套路。验收应优先依赖外部 outcome、可执行检查和反事实测试，而不是只奖励“看起来像推理”。
\end{warningbox}

\lecturefigure{slide-13-create-build-make.jpg}{Create, build, and make：能力扩张必须落到可完成的作品}{00:09:03--00:09:35}

\teachervoice{00:07:58--00:09:34，课堂把 next-token prediction 与 RL on CoT 称为两种 scaling paradigm。老师强调的是监督信号发生变化：前者学习分布，后者可把复杂任务的结果反馈写进训练；这并不意味着所有长任务都已被解决。}

\subsection{两条路径：Capability-first 与 Product-belief-first}

第一条路径从新能力出发，寻找用户熟悉的 form factor；第二条路径先提出产品信念，再构造数据、eval 与环境让模型呈现该行为。前者适合探索“模型突然会了什么”，后者适合追求长期一致的产品体验。真实项目通常在两者之间往返。

\lecturefigure{slide-14-two-product-research-paths.jpg}{Research-driven product 的两条路径}{00:09:34--00:10:00}

\lecturefigure{slide-15-familiar-form-unfamiliar-capability.jpg}{路径一：用熟悉 form factor 承载陌生能力}{00:09:43--00:10:06}

\begin{importantbox}{Capability、Affordance 与 Eval 的三角形}
\term{Capability} 是模型潜在能做什么；\term{affordance} 是界面允许用户怎样调用、编辑和撤销；\term{eval} 是系统用什么证据判断完成。产品失败常来自三者错位：模型会做但界面不能表达，界面能触发但 eval 不测真实结果，或 eval 优化了与用户目标不同的代理指标。
\end{importantbox}

\subsection{案例一：CLIP fashion search}

CLIP 把图像和文本映射到共享 embedding 空间。设图像编码为 $f_I(i)$、文本编码为 $f_T(q)$，检索得分常用余弦相似度
\[
s(i,q)=\frac{f_I(i)^{\top}f_T(q)}{\lVert f_I(i)\rVert\lVert f_T(q)\rVert}.
\]
这项陌生能力被放入熟悉的搜索界面：用户给出文本或目标服装，系统返回视觉相似结果。读三张连续状态时，应比较 query 意图怎样改变结果集，而不是把任何一组漂亮图片都当作搜索成功。

\lecturefigure{slide-16-clip-fashion-target.jpg}{Fashion search：从目标服装与视觉属性开始检索}{00:10:06--00:10:11}

\lecturefigure{slide-17-clip-fashion-yellow-search.jpg}{Fashion search：黄色与款式语义改变结果分布}{00:10:11--00:10:23}

\lecturefigure{slide-18-clip-fashion-style-search.jpg}{Fashion search：自然语言风格查询返回新的候选集合}{00:10:23--00:10:54}

\begin{knowledgebox}{读图：检索 demo 至少要测四件事}
第一是 relevance，结果是否匹配 query；第二是 diversity，是否只返回近重复；第三是 controllability，修改颜色、材质或轮廓能否局部改变结果；第四是 failure visibility，系统是否让用户看出为什么匹配。共享 embedding 提供排序机制，但不自动保证公平、库存可用或风格词在不同文化中的一致含义。
\end{knowledgebox}

\subsection{案例二：100K context 与文件上传}

长上下文能力进入产品后，最直观 form factor 是“上传文件并提问”。然而 token window 足够大，只说明输入可被接收，不说明模型找到了关键页、正确聚合数字或保持来源对应。一个 85 页 10-K 分析任务应至少验收 retrieval、calculation、citation 与 omission 四类错误。

\lecturefigure{slide-19-100k-context-business-analyst.jpg}{100K context 的产品承诺：让 Claude 充当企业分析师}{00:10:54--00:11:00}

\lecturefigure{slide-20-100k-context-input.jpg}{输入侧：85 页 Form 10-K 与具体分析请求}{00:11:00--00:11:13}

\lecturefigure{slide-21-100k-context-analysis.jpg}{输出侧：关键财务指标、表格与解释被组织为分析}{00:11:13--00:11:34}

\begin{warningbox}{Long context 不等于 long-context reliability}
上下文窗口是容量上限；可靠性还取决于位置偏差、跨页计算、表格解析、冲突证据和引用。若产品只测“能否上传”，模型可能在最危险的情况下给出流畅但遗漏关键负债的摘要。应构造可定位答案、跨段聚合、无答案和冲突文件四组测试。
\end{warningbox}

\subsection{案例三：Summarizer 与不确定性界面}

Summarizer demo 把重复段落、校准和编辑集中到文档界面。它表达的产品信念不是“摘要越短越好”，而是让用户看见哪些句子被保留、哪些内容重复、哪些结论需要再核查。若系统能产生 $K$ 个候选，可把候选一致性当作不确定性线索，但不能把一致误差当作置信。与前面的 100K context 案例相比，这里更关心用户怎样审查和修改结果：长上下文解决“能读多少”，summarizer surface 解决“怎样把证据压缩成仍可追踪的 artifact”。

\lecturefigure{slide-22-o1-summarizer-calibration.jpg}{Summarizer：重复检测、P(IK) 与可编辑文档界面}{00:12:41--00:13:39}

\begin{knowledgebox}{从产品界面到评测样本}
一个实际工作流可以把用户的接受、修改、删除、重试和引用跳转记录成结构化 evidence。最小事件可写为
\[
e=(\text{task},\text{model version},\text{action},\text{artifact diff},\text{outcome}).
\]
只有 outcome 被定义，日志才可能转成训练或 eval 数据；单纯收集“用户点了哪里”容易把困惑误判成偏好。
\end{knowledgebox}

\begin{lstlisting}[caption={Capability-first 原型的最小验收循环}]
task = observe_real_user_job()
prototype = expose_capability_in_familiar_surface(task)
trace = run_with_permissions_and_logging(prototype)
failures = classify(trace, by=[outcome, control, trust, cost])
eval_set = convert_failures_to_reproducible_cases(failures)
ship_only_if(eval_set, product_constraints).passes()
\end{lstlisting}

\teachervoice{00:09:34--00:13:39，讲者用 Inter Alia、100K context 和 summarizer 说明第一条路径：先出现陌生能力，再寻找熟悉 form factor。实践经验是界面必须让用户理解、控制并验证能力，而不是仅把模型接到输入框。}

\subsection{本章小结}

Next-token prediction 提供广泛先验，RL on CoT 可以把轨迹 outcome 写进目标；二者都不能单独决定产品。Capability-first 路径需要合适 form factor，product-belief-first 路径需要把信念翻译成训练与评测。CLIP、长上下文和 summarizer 的共同教训是：产品层必须显式定义可控性与证据。

